% Problem Set 8 -- Part 1: questions (a), (b), (c). Numbers from ps8_partA.py.

\section{Questions (a)--(c): Seasonality in Returns}

All regressions use the 75 portfolios sorted on size and past 1-month, 2--12-month and
13--60-month returns. The left-hand side is the portfolio return in excess of the T-bill, in
\% per month. The sample is 1931:02--2026:06 ($T=1{,}145$), the months in which all three
sorts exist. Four portfolios have a few months coded $-99.99$, which we drop, so their $T$ is
slightly smaller (at least 1{,}127).

\subsection*{Question (a): the January seasonal}

We run $r^e_{i,t} = c_i + J_i\,\mathrm{Jan}_t + \varepsilon_{i,t}$ for each portfolio. The
intercept $c_i$ is the mean excess return from February to December, and $J_i$ is the extra
January return. Table~\ref{tab:jan} reports $J_i$ for all 75 portfolios.

\begin{table}[H]
\centering\small
\caption{January premium $J_i$ (\% per month) from $r^e_{i,t}=c_i+J_i\,\mathrm{Jan}_t+\varepsilon_{i,t}$, OLS $t$-statistics in parentheses. Rows are size quintiles, columns are past-return quintiles (Low = past losers, High = past winners). 1931:02--2026:06.}
\label{tab:jan}
\begin{tabular}{lccccc}
\toprule
Size & Low & 2 & 3 & 4 & High \\
\midrule
\multicolumn{6}{l}{\textit{Panel A: sorted on past 1-month return}} \\
Small & 7.32 (6.5) & 4.74 (4.8) & 4.15 (4.3) & 3.72 (4.1) & 3.49 (3.7) \\
2 & 3.92 (4.0) & 3.05 (3.7) & 2.33 (3.0) & 2.14 (2.6) & 1.32 (1.6) \\
3 & 3.02 (3.4) & 2.35 (3.1) & 1.78 (2.5) & 1.44 (2.1) & 0.39 (0.5) \\
4 & 2.15 (2.6) & 1.69 (2.4) & 1.21 (1.9) & 0.40 (0.6) & $-$0.10 ($-$0.1) \\
Big & 1.16 (1.6) & 0.54 (0.9) & 0.15 (0.3) & $-$0.11 ($-$0.2) & $-$0.52 ($-$0.8) \\
\midrule
\multicolumn{6}{l}{\textit{Panel B: sorted on past 2--12-month return}} \\
Small & 7.45 (6.7) & 5.14 (5.3) & 4.32 (4.8) & 3.56 (3.8) & 3.57 (3.9) \\
2 & 4.07 (3.9) & 3.10 (3.6) & 2.38 (3.1) & 1.68 (2.2) & 1.66 (2.0) \\
3 & 2.95 (3.0) & 2.10 (2.6) & 1.68 (2.3) & 1.15 (1.7) & 0.91 (1.3) \\
4 & 2.23 (2.2) & 1.45 (1.9) & 0.95 (1.4) & 0.39 (0.6) & 0.34 (0.5) \\
Big & 1.69 (1.9) & 0.82 (1.2) & 0.13 (0.2) & $-$0.32 ($-$0.6) & $-$0.55 ($-$0.9) \\
\midrule
\multicolumn{6}{l}{\textit{Panel C: sorted on past 13--60-month return}} \\
Small & 7.42 (6.9) & 4.73 (5.1) & 3.45 (4.0) & 4.24 (4.6) & 3.70 (3.7) \\
2 & 4.56 (4.6) & 2.90 (3.4) & 2.08 (2.7) & 1.71 (2.1) & 1.61 (1.9) \\
3 & 3.82 (4.1) & 2.17 (2.8) & 1.73 (2.4) & 1.23 (1.8) & 0.99 (1.3) \\
4 & 2.67 (3.0) & 1.29 (1.7) & 1.20 (1.7) & 0.77 (1.1) & 0.28 (0.4) \\
Big & 2.06 (3.1) & 0.72 (1.2) & 0.31 (0.5) & $-$0.15 ($-$0.3) & $-$0.34 ($-$0.6) \\
\bottomrule
\end{tabular}
\end{table}

\begin{tight}
  \item \textbf{Is there a strong January seasonal different from other months?} Yes, January
  is by far the strongest month for most of these portfolios. The January premium is positive
  and significant ($t>2$) in 43 of the 75 portfolios, and across all 75 it averages 2.1\% on
  top of a February--December mean of about 0.8\%. In other words, an average portfolio earns
  almost four times its normal monthly return in January.

  \item \textbf{For which portfolios is this largest or smallest?} The effect is largest for
  small past losers and disappears for large past winners. In every panel the January premium
  generally falls as we move from small to big stocks and from losers to winners, so the small--loser
  corner earns more than 7\% extra in January with $t$-statistics near 7. The big--winner
  corner has a small negative premium that is not significant. The pattern looks the same
  whether ``losers'' are defined over 1, 12 or 60 months, so the January effect is a
  small-stock and past-loser effect at every horizon.

  \item \textbf{Why might this be?} The most common explanation is tax-loss selling. Investors
  sell stocks with capital losses in December to cut their tax bill, which pushes down prices
  of past losers, and then buy back in January when the selling pressure ends. Small stocks are
  hit hardest because they are illiquid, held mostly by taxable individuals, and expensive to
  arbitrage. Window dressing adds to this: fund managers dump embarrassing losers before
  year-end reports and buy riskier small stocks again in January. The intercepts also show that
  past 12-month losers earn only 0.1--0.3\% per month from February to December against
  1.0--1.3\% for winners, so momentum is a February--December phenomenon that reverses in
  January. All of
  these stories point to year-end trading pressure, not a change in risk, as the source of the
  January return.
\end{tight}

\subsection*{Question (b): adding December}

We add a December dummy,
$r^e_{i,t} = c_i + J_i\,\mathrm{Jan}_t + D_i\,\mathrm{Dec}_t + \varepsilon_{i,t}$, so $c_i$ is
now the February--November mean. Table~\ref{tab:dec} reports $D_i$ for all 75 portfolios, and
Table~\ref{tab:corners} compares the three means in the four corner portfolios.

\begin{table}[H]
\centering\small
\caption{December premium $D_i$ (\% per month) from $r^e_{i,t}=c_i+J_i\,\mathrm{Jan}_t+D_i\,\mathrm{Dec}_t+\varepsilon_{i,t}$, OLS $t$-statistics in parentheses. Layout as in Table~\ref{tab:jan}.}
\label{tab:dec}
\begin{tabular}{lccccc}
\toprule
Size & Low & 2 & 3 & 4 & High \\
\midrule
\multicolumn{6}{l}{\textit{Panel A: sorted on past 1-month return}} \\
Small & $-$0.12 ($-$0.1) & 0.10 (0.1) & $-$0.06 ($-$0.1) & $-$0.01 (0.0) & 0.29 (0.3) \\
2 & 0.68 (0.7) & 0.76 (0.9) & 0.74 (0.9) & 0.99 (1.2) & 1.11 (1.3) \\
3 & 1.02 (1.1) & 0.69 (0.9) & 1.12 (1.6) & 1.02 (1.5) & 1.03 (1.3) \\
4 & 0.94 (1.1) & 0.80 (1.2) & 1.16 (1.8) & 0.94 (1.4) & 1.17 (1.6) \\
Big & 0.43 (0.6) & 0.52 (0.9) & 0.67 (1.2) & 0.78 (1.4) & 0.60 (0.9) \\
\midrule
\multicolumn{6}{l}{\textit{Panel B: sorted on past 2--12-month return}} \\
Small & $-$0.65 ($-$0.6) & $-$0.32 ($-$0.3) & $-$0.24 ($-$0.3) & 0.30 (0.3) & 1.00 (1.1) \\
2 & 0.16 (0.2) & 0.35 (0.4) & 0.79 (1.0) & 1.39 (1.8) & 1.46 (1.8) \\
3 & 0.62 (0.6) & 0.66 (0.8) & 0.68 (0.9) & 1.27 (1.9) & 1.48 (2.0) \\
4 & 0.34 (0.3) & 0.61 (0.8) & 0.79 (1.2) & 1.13 (1.8) & 1.60 (2.4) \\
Big & $-$0.60 ($-$0.7) & $-$0.01 (0.0) & 0.52 (0.9) & 0.90 (1.6) & 1.19 (2.0) \\
\midrule
\multicolumn{6}{l}{\textit{Panel C: sorted on past 13--60-month return}} \\
Small & $-$0.31 ($-$0.3) & 0.07 (0.1) & 0.22 (0.2) & $-$0.10 ($-$0.1) & 0.07 (0.1) \\
2 & 0.58 (0.6) & 0.91 (1.1) & 0.94 (1.2) & 0.75 (0.9) & 0.51 (0.6) \\
3 & 1.00 (1.1) & 0.92 (1.2) & 1.10 (1.5) & 0.71 (1.0) & 0.74 (1.0) \\
4 & 0.66 (0.7) & 0.80 (1.1) & 0.98 (1.4) & 1.03 (1.5) & 0.87 (1.2) \\
Big & 0.65 (1.0) & 0.68 (1.1) & 0.54 (0.9) & 0.67 (1.1) & 0.39 (0.6) \\
\bottomrule
\end{tabular}
\end{table}

\begin{table}[H]
\centering\small
\caption{Corner portfolios from the January--December regression. $c$ is the February--November mean, $J$ and $D$ are the extra January and December returns (\% per month, $t$ in parentheses), and the last column is the $t$-statistic of $J-D$.}
\label{tab:corners}
\begin{tabular}{llcccc}
\toprule
Sort & Portfolio & Feb--Nov $c$ & January $J$ & December $D$ & $t(J-D)$ \\
\midrule
1-month & Small, low & 1.42 & 7.31 (6.5) & $-$0.12 ($-$0.1) & 4.9 \\
 & Small, high & $-$0.23 & 3.51 (3.7) & 0.29 (0.3) & 2.5 \\
 & Big, low & 0.70 & 1.20 (1.6) & 0.43 (0.6) & 0.8 \\
 & Big, high & 0.47 & $-$0.46 ($-$0.7) & 0.60 (0.9) & $-$1.2 \\
\addlinespace
2--12-month & Small, low & 0.12 & 7.39 (6.6) & $-$0.65 ($-$0.6) & 5.3 \\
 & Small, high & 1.15 & 3.66 (4.0) & 1.00 (1.1) & 2.2 \\
 & Big, low & 0.36 & 1.64 (1.8) & $-$0.60 ($-$0.7) & 1.8 \\
 & Big, high & 0.89 & $-$0.44 ($-$0.7) & 1.19 (2.0) & $-$2.0 \\
\addlinespace
13--60-month & Small, low & 0.85 & 7.39 (6.8) & $-$0.31 ($-$0.3) & 5.3 \\
 & Small, high & 0.50 & 3.70 (3.7) & 0.07 (0.1) & 2.7 \\
 & Big, low & 0.57 & 2.12 (3.2) & 0.65 (1.0) & 1.6 \\
 & Big, high & 0.63 & $-$0.31 ($-$0.5) & 0.39 (0.6) & $-$0.8 \\
\bottomrule
\end{tabular}
\end{table}

\begin{tight}
  \item \textbf{Is there a strong December and January seasonal, and are they different from
  each other?} Yes for January, but December is only a modest seasonal, and the two are
  clearly different for small stocks. Adding the December dummy barely changes the January
  premia. December premia are mostly positive, averaging 0.6\% across the 75 portfolios, but
  only 2 of 75 have $t>2$ with OLS standard errors. December is a calmer month than the rest of
  the year, so with White standard errors 12 portfolios reach $t>2$, 8 of them 12-month winners.
  December is therefore a real but modest seasonal, about 1--1.6\% in a few portfolios. January exceeds December significantly in 21 portfolios, all but one of them in
  the two smallest size quintiles, while for big stocks the two months cannot be told apart.

  \item \textbf{For which portfolios are these seasonalities largest or smallest?} The two
  seasonals run in opposite directions across portfolios. The January effect is largest for
  small losers, while December is weakest there: small stocks earn nothing significant extra in December,
  and small 12-month losers are even slightly negative. December is largest for mid-to-large
  12-month winners, where it is about 1--1.6\% and significant with White standard errors, and these are
  exactly the portfolios with no significant January effect. So small losers do badly into year-end and then rebound,
  while large winners do well into year-end and, if anything, give a little back in January.

  \item \textbf{Can you conjecture why these differences exist?} The mirror-image pattern is
  what year-end selling pressure would produce. In December, tax-loss sellers and window
  dressers sell losers and hold or buy winners, so losers lag and winners gain. In January the
  pressure lifts, and the losers bounce back while winners are sold to realize gains in the new
  tax year. The effect concentrates in small stocks because their prices move most for a given
  amount of selling. The December premium in mid and large stocks may also reflect the
  pre-holiday effect and year-end inflows into the broad market, which would explain why mid
  and large stocks earn a little extra in December across most past-return groups.
\end{tight}

\subsection*{Question (c): market efficiency}

\begin{tight}
  \item \textbf{Are these findings evidence for or against market efficiency?} At face value
  they are evidence against it, but weaker evidence than the $t$-statistics suggest. A return
  pattern tied to the calendar is known in advance, so in an efficient market investors should
  buy small losers in late December and trade the January premium away. Its strength over 95 years, at
  $t$-statistics near 7, is hard to explain with risk. Still, several caveats
  below mean the findings do not cleanly reject efficiency.

  \item \textbf{How would you try to reconcile them with the efficient markets hypothesis?}
  There are three ways. First, statistics: we looked at 75 portfolios and could have tested
  any of 12 months, so some large $t$-statistics are expected by chance. The effect has also
  weakened sharply since it was published in the late 1970s and early 1980s: the January
  premium is significant in 46 of 75 portfolios before 1983 but in only 6 after, which is what
  we would expect if investors learned about it and traded it away. Second, transaction costs: small losers
  have wide bid--ask spreads and low prices, so the premium may not survive realistic trading
  costs and short-sale limits. Third, the model: we assumed expected returns are constant across
  months, and EMH only says prices reflect information given the right model of risk. If small
  stocks carry more risk in January, or if taxes make the after-tax return equal across months,
  then a higher pre-tax January return is just compensation. This is the joint-hypothesis
  problem: we are always testing efficiency and the risk model together.

  \item \textbf{How are the findings perhaps a violation of efficient markets?} The risk story
  is hard to defend, because nothing about a firm's risk changes on January 1. The tax story
  only explains why some investors sell, not why arbitrageurs without tax motives fail to buy.
  The effect is also very large relative to normal returns and has been documented in many
  countries, including ones with different tax years (Gultekin and Gultekin, 1983), which makes pure data mining less
  likely. If the premium survives after costs, it is a violation at least of the weak form,
  since the signal uses only past prices and the calendar. Either way, it shows that limits to
  arbitrage let predictable patterns persist in the parts of the market that are hardest to
  trade.
\end{tight}
